\subsection{The Namespace Graph: Definitions and Examples}
\label{sec:namespace-graph}

This appendix collects the formal definitions, worked examples, and illustrative figures for the namespace graphs. Extended statistical analysis (degree distributions, centrality, community detection, robustness curves) appears in Appendix~\ref{app:ns-detail}.

\subsubsection{Namespaces versus modules}
Each declaration carries a fully qualified name (e.g., \decl{Nat.Prime.dvd\_mul}) whose dot-separated components form a hierarchy determined by \texttt{namespace} blocks in the source code, not by the file path. The following example illustrates how Lean's naming mechanisms determine fully qualified names:

\begin{codebox}
\textit{\color{gray}-- \textbf{Namespace blocks} determine the declaration's qualified name:}\\
\textbf{namespace} \textcolor{green!50!black}{Nat}\\
\hspace{1.5em}\textbf{theorem} \textcolor{red!70!black}{add\_comm} (a b : Nat) : a + b = b + a := \ldots\\
\hspace{1.5em}\textit{\color{gray}-- fully qualified name: \textcolor{red!70!black}{Nat.add\_comm} \quad(depth 1)}\\[3pt]
\hspace{1.5em}\textbf{namespace} \textcolor{green!50!black}{Prime}\\
\hspace{3em}\textbf{theorem} \textcolor{red!70!black}{dvd\_mul} (hp : Prime p) : p | a * b $\to$ \ldots\ := \ldots\\
\hspace{3em}\textit{\color{gray}-- fully qualified name: \textcolor{red!70!black}{Nat.Prime.dvd\_mul} \quad(depth 2)}\\
\hspace{1.5em}\textbf{end} \textcolor{green!50!black}{Prime}\\
\textbf{end} \textcolor{green!50!black}{Nat}\\[6pt]
\textit{\color{gray}-- \textbf{File path} determines the module name (a separate hierarchy):}\\
\textit{\color{gray}-- File:}\quad \textcolor{blue!60!black}{Mathlib/Data/Nat/Prime/Basic.lean}\\
\textit{\color{gray}-- Module:}\quad \textcolor{blue!60!black}{Mathlib.Data.Nat.Prime.Basic}
\end{codebox}

\noindent
Each declaration's fully qualified name is a dot-separated sequence determined by its enclosing \texttt{namespace} blocks (\textcolor{green!50!black}{green}), not by the file path. The file path determines the \emph{module} name (\textcolor{blue!60!black}{blue}), which governs \texttt{import} visibility but does not affect qualified names. In practice, \module{Mathlib} conventions strongly correlate the two hierarchies---declarations in \ns{Nat.Prime} typically live in modules under \module{Mathlib.Data.Nat.Prime}---but the correlation is a social convention, not a language rule. A single namespace may span many files, and a single file may contribute declarations to multiple namespaces (\S\ref{sec:ns-mod-cross}).

Truncating names to varying depths yields an intermediate family of graphs between the coarse module level and the fine declaration level.

The $308{,}129$ declarations in \module{Mathlib} span six depth levels. Table~\ref{tab:ns-depth-count} shows how the number of distinct namespaces grows with truncation depth.
\begin{table}[ht]
\centering
\caption{Number of distinct namespaces at each truncation depth~$k$.}
\label{tab:ns-depth-count}
\begin{tabular}{cr}
\toprule
Depth $k$ & $|\mathcal{N}_k|$ \\
\midrule
1          & $3{,}184$ \\
2          & $10{,}097$ \\
3          & $13{,}785$ \\
4          & $15{,}198$ \\
5          & $15{,}443$ \\
6 (max)    & $15{,}456$ \\
\bottomrule
\end{tabular}
\end{table}

\noindent
Growth saturates rapidly: $90\%$ of all namespaces are distinguished by depth~$3$. Declarations with names of the form $X.y$ (a single namespace component plus a short name) account for $67\%$ of all declarations; only $10\%$ have depth $\ge 4$. The naming convention of \module{Mathlib} is, in practice, shallow.

\begin{definition}[Namespace graph at depth~$k$]
\label{def:ns-graph}
For each depth $k \ge 1$, let $\mathrm{ns}_k(d)$ denote the namespace of declaration~$d$ truncated to depth~$k$: the first $k$ dot-separated components of $d$'s fully qualified name. If $d$'s name has fewer than $k+1$ components, $\mathrm{ns}_k(d)$ defaults to the full parent namespace; if $d$'s name has no dot, $\mathrm{ns}_k(d) = \ns{\_root\_}$. The \emph{namespace graph at depth~$k$} is the weighted directed graph
\[
  G_{\mathrm{ns}}^{(k)} = (\mathcal{N}_k,\, E_k),
\]
where $\mathcal{N}_k = \{\mathrm{ns}_k(d) \mid d \in \mathcal{D}\}$ is the set of all distinct depth-$k$ namespaces, and each edge $(n_1, n_2) \in E_k$ carries a weight equal to the number of edges $(d_1, d_2) \in E_{\mathrm{thm}}$ such that $\mathrm{ns}_k(d_1) = n_1 \neq n_2 = \mathrm{ns}_k(d_2)$.
\end{definition}

\noindent At $k=1$, the $3{,}184$ coarse namespaces correspond roughly to mathematical topics; at $k=6$, the $15{,}456$ fine-grained namespaces approach the resolution of individual naming scopes. The aggregation process is illustrated in Figure~\ref{fig:ns-aggregation}.

\begin{figure}[!htb]
\centering
\begin{tikzpicture}[
  every node/.style={font=\scriptsize, inner sep=2pt},
  dot/.style={circle, fill=red!70!black, minimum size=3pt, inner sep=0pt},
  dep/.style={->, >=Stealth, red!70!black, semithick},
  depx/.style={->, >=Stealth, red!70!black, semithick, dashed},
  nsbox/.style={rounded corners=5pt, draw=green!50!black, thick, fill=green!10},
  agg/.style={->, >=Stealth, gray!50, very thick, dashed, shorten <=2pt, shorten >=2pt},
]
  \fill[blue!6, rounded corners=3pt] (-0.3,-2.1) rectangle (2.5,1.8);
  \fill[blue!6, rounded corners=3pt] (3.5,-0.6) rectangle (5.6,1.8);
  \fill[green!8, rounded corners=3pt] (0.0,-0.15) rectangle (5.3,1.25);
  \fill[green!8, rounded corners=2pt] (0.0,-1.75) rectangle (2.2,-0.55);
  \draw[blue!50!black, rounded corners=3pt, thick, dashed]
    (-0.3,-2.1) rectangle (2.5,1.8);
  \node[font=\tiny\sffamily\bfseries, blue!60!black, anchor=north west]
    at (-0.2,1.75) {\textnormal{\module{Data.Nat.Basic}}};
  \draw[blue!50!black, rounded corners=3pt, thick, dashed]
    (3.5,-0.6) rectangle (5.6,1.8);
  \node[font=\tiny\sffamily\bfseries, blue!60!black, anchor=north west]
    at (3.6,1.75) {\textnormal{\module{Data.Nat.Defs}}};
  \draw[green!50!black, rounded corners=3pt, thick, dashed]
    (0.0,-0.15) rectangle (5.3,1.25);
  \node[font=\tiny\sffamily\bfseries, green!40!black, anchor=south west]
    at (0.1,1.3) {\ns{Nat}};
  \draw[green!50!black, rounded corners=2pt, thick, dashed]
    (0.0,-1.75) rectangle (2.2,-0.55);
  \node[font=\tiny\sffamily\bfseries, green!40!black, anchor=south west]
    at (0.1,-0.5) {\ns{Nat.Primrec}};
  \node[dot, label={[font=\tiny\ttfamily, red!70!black]left:add\_comm}]
    (ac) at (1.8,0.9) {};
  \node[dot, label={[font=\tiny\ttfamily, red!70!black]left:add\_left\_comm}]
    (alc) at (0.5,0.15) {};
  \node[dot, label={[font=\tiny\ttfamily, red!70!black]right:Primrec.zero}]
    (prz) at (0.6,-1.15) {};
  \node[dot, label={[font=\tiny\ttfamily, red!70!black]right:succ\_eq\_one}]
    (sea) at (4.1,0.8) {};
  \node[dot, label={[font=\tiny\ttfamily, red!70!black]right:zero\_add}]
    (za) at (4.1,0.2) {};
  \draw[dep] (alc) -- (ac);
  \draw[dep, dashed] (sea) to[out=175,in=10] (ac);
  \draw[dep, dashed] (za) to[out=175,in=-10] (alc);
  \draw[dep, densely dotted, thick] (prz) to[out=70,in=-100] (ac);
  \draw[->, >=Stealth, blue!60!black, thick]
    (3.5,1.55) -- node[above, font=\tiny\sffamily, text=blue!60!black] {\texttt{import}} (2.5,1.55);
  \node[font=\small\sffamily, anchor=north] at (2.6,-2.5) {(a) $G_{\mathrm{thm}}$};
  \draw[agg] (5.9,-0.1) -- node[above, font=\tiny\sffamily, text=gray!60] {aggregate}
    node[below, font=\tiny\sffamily, text=gray!60] {depth $2$} (7.3,-0.1);
  \node[nsbox, minimum width=2.2cm, minimum height=1.1cm]
    (nat2) at (8.8,0.55) {};
  \node[font=\small\sffamily\bfseries, green!50!black] at (8.8,0.75) {\ns{Nat}};
  \node[font=\tiny\sffamily, text=gray] at (8.8,0.3) {4 declarations};
  \node[nsbox, minimum width=2.2cm, minimum height=0.85cm]
    (npr2) at (8.8,-1.15) {};
  \node[font=\scriptsize\sffamily\bfseries, green!50!black] at (8.8,-0.97) {\ns{Nat.Primrec}};
  \node[font=\tiny\sffamily, text=gray] at (8.8,-1.35) {1 declaration};
  \draw[->, >=Stealth, green!50!black, thick] (npr2) --
    node[right, font=\tiny, text=green!50!black] {$w\!=\!1$} (nat2);
  \node[font=\tiny\itshape, text=gray, align=center] at (8.8,-2.0)
    {\decl{Primrec.zero} $\to$ \decl{add\_comm}};
  \draw[green!50!black, ->, >=Stealth, thin, opacity=0.35]
    ([xshift=-4pt]nat2.north) to[out=130,in=50,looseness=6]
    ([xshift=4pt]nat2.north);
  \node[font=\tiny, text=green!40!black] at (8.8,1.75) {$3$ internal};
  \node[font=\small\sffamily, anchor=north] at (8.8,-2.5) {(b) $G_{\mathrm{ns}}^{(2)}$};
  \draw[agg] (10.3,-0.1) -- node[above, font=\tiny\sffamily, text=gray!60] {aggregate}
    node[below, font=\tiny\sffamily, text=gray!60] {depth $1$} (11.7,-0.1);
  \node[nsbox, minimum width=2.2cm, minimum height=2.2cm]
    (nat1) at (13.0,-0.1) {};
  \node[font=\small\sffamily\bfseries, green!50!black] at (13.0,0.4) {\ns{Nat}};
  \node[font=\tiny\sffamily, text=gray] at (13.0,0.05) {5 declarations};
  \node[font=\tiny\sffamily, text=gray] at (13.0,-0.3) {4 internal edges};
  \node[font=\tiny\itshape, text=red!50!black] at (13.0,-0.7) {contain.\,$=100\%$};
  \node[font=\small\sffamily, anchor=north] at (13.0,-2.5) {(c) $G_{\mathrm{ns}}^{(1)}$};
\end{tikzpicture}
\caption{Aggregation from $G_{\mathrm{thm}}$ to $G_{\mathrm{ns}}^{(k)}$. (a)~Declarations in~$G_{\mathrm{thm}}$; \textcolor{green!40!black}{green} = namespace boundaries, \textcolor{blue!60!black}{blue} = module boundaries. (b)~Depth-$2$ namespaces. (c)~Depth~$1$. All names have the \texttt{Nat.}\ prefix removed.}
\label{fig:ns-aggregation}
\label{fig:ns-aggregation-app}
\end{figure}

\noindent\textbf{Cycles in the namespace graph.}
\label{rem:ns-cycles}
Unlike $G_{\mathrm{module}}$ and $G_{\mathrm{thm}}$, which are directed acyclic graphs---the compiler and the type-checker respectively forbid circular dependencies---the namespace graph $G_{\mathrm{ns}}^{(k)}$ can contain directed cycles. Two namespaces may depend on each other: namespace~$A$ contains a declaration that cites a declaration in namespace~$B$, while $B$ contains a different declaration that cites one in~$A$. The underlying declaration-level graph remains acyclic, but the aggregation into namespaces collapses the strict topological ordering.

At depth~$2$, $G_{\mathrm{ns}}^{(2)}$ contains $38$ strongly connected components with more than one node, the largest comprising $5{,}899$ namespaces---$58\%$ of all depth-$2$ namespaces. This giant component reflects the deeply interconnected core of Mathlib: most mathematical sub-disciplines at this granularity depend mutually on one another. This structural difference---DAGs at the module and declaration levels, cyclic graphs at the namespace level---reflects the fact that namespace boundaries, unlike file boundaries, are not constrained by a compiler-enforced import ordering.

\begin{figure}[ht]
\centering
\begin{tikzpicture}[
  every node/.style={font=\scriptsize, inner sep=2pt},
  dot/.style={circle, fill=red!70!black, minimum size=3pt, inner sep=0pt},
  dep/.style={->, >=Stealth, red!70!black, semithick},
  nsbox/.style={rounded corners=5pt, draw=green!50!black, thick, fill=green!10,
                minimum width=2.8cm, minimum height=1.6cm},
  agg/.style={->, >=Stealth, gray!50, very thick, dashed,
              shorten <=2pt, shorten >=2pt},
]
  % ============ Panel (a): Declaration level ============
  \node[font=\small\sffamily, anchor=south] at (2.0,2.5) {(a) $G_{\mathrm{thm}}$ (acyclic)};
  % Namespace boxes
  \fill[green!8, rounded corners=3pt] (-0.3,-0.3) rectangle (2.0,2.2);
  \draw[green!50!black, rounded corners=3pt, thick, dashed]
    (-0.3,-0.3) rectangle (2.0,2.2);
  \node[font=\tiny\sffamily\bfseries, green!40!black, anchor=north west]
    at (-0.2,2.15) {\ns{Nat}};

  \fill[green!8, rounded corners=3pt] (2.5,-0.3) rectangle (4.8,2.2);
  \draw[green!50!black, rounded corners=3pt, thick, dashed]
    (2.5,-0.3) rectangle (4.8,2.2);
  \node[font=\tiny\sffamily\bfseries, green!40!black, anchor=north west]
    at (2.6,2.15) {\ns{Int}};

  % Declarations
  \node[dot, label={[font=\tiny\ttfamily, red!70!black, align=center]below:prime\_iff\_\\[-1pt]prime\_int}]
    (npip) at (1.0,1.4) {};
  \node[dot, label={[font=\tiny\ttfamily, red!70!black]below:lcm}]
    (nlcm) at (1.0,0.3) {};

  \node[dot, label={[font=\tiny\ttfamily, red!70!black]below:dvd\_natAbs}]
    (idna) at (3.5,1.4) {};
  \node[dot, label={[font=\tiny\ttfamily, red!70!black]below:lcm\_neg}]
    (ilcm) at (3.5,0.3) {};

  % Declaration-level edges (no cycle!)
  \draw[dep] (npip) -- (idna);
  \draw[dep] (ilcm) -- (nlcm);

  % ============ Arrow a → b ============
  \draw[agg] (5.3,0.95) -- node[above, font=\tiny\sffamily, text=gray!60] {aggregate}
    (7.0,0.95);

  % ============ Panel (b): Namespace level ============
  \node[font=\small\sffamily, anchor=south] at (9.5,2.5) {(b) $G_{\mathrm{ns}}^{(2)}$ (cyclic)};
  \node[nsbox, minimum width=1.6cm, minimum height=1.3cm] (ns_nat) at (8.0,0.95) {};
  \node[font=\scriptsize\sffamily\bfseries, green!50!black] at (8.0,1.1) {\ns{Nat}};
  \node[font=\tiny\sffamily, text=gray] at (8.0,0.7) {2 decls};

  \node[nsbox, minimum width=1.6cm, minimum height=1.3cm] (ns_int) at (11.0,0.95) {};
  \node[font=\scriptsize\sffamily\bfseries, green!50!black] at (11.0,1.1) {\ns{Int}};
  \node[font=\tiny\sffamily, text=gray] at (11.0,0.7) {2 decls};

  % Namespace-level edges: cycle!
  \draw[->, >=Stealth, green!50!black, thick]
    ([yshift=2pt]ns_nat.north east) to[out=30,in=150]
    node[above, font=\tiny, text=green!40!black] {$w\!=\!1$}
    ([yshift=2pt]ns_int.north west);
  \draw[->, >=Stealth, green!50!black, thick]
    ([yshift=-2pt]ns_int.south west) to[out=210,in=330]
    node[below, font=\tiny, text=green!40!black] {$w\!=\!1$}
    ([yshift=-2pt]ns_nat.south east);
\end{tikzpicture}
\caption{How aggregation creates cycles. (a)~At the declaration level, two edges cross the \ns{Nat}--\ns{Int} boundary in opposite directions: \decl{Nat.prime\_iff\_prime\_int} cites \decl{Int.dvd\_natAbs}, and \decl{Int.lcm\_neg} cites \decl{Nat.lcm}. No cycle exists among these four declarations. (b)~After aggregation to depth-$2$ namespaces, both edges collapse to namespace-level edges, creating a directed cycle \ns{Nat} $\to$ \ns{Int} $\to$ \ns{Nat}. The declaration-level DAG property is lost.}
\label{fig:ns-cycle-example}
\end{figure}

\paragraph{Containment Decay}
\label{sec:containment-decay}

We now quantify how namespace containment erodes as the grouping becomes finer.

\begin{definition}[Containment ratio]
\label{def:containment}
The \emph{containment ratio at depth~$k$} is the proportion of edges in $G_{\mathrm{thm}}$ whose endpoints share the same depth-$k$ namespace:
\[
  \mathrm{contain}(k) \;=\; \frac{|E_{\mathrm{intra}}^{(k)}|}{|E_{\mathrm{thm}}|},
  \qquad
  E_{\mathrm{intra}}^{(k)} = \{(d_1,d_2) \in E_{\mathrm{thm}} \mid \mathrm{ns}_k(d_1) = \mathrm{ns}_k(d_2)\}.
\]
\end{definition}

\noindent
This definition parallels the module containment ratio (Definition~\ref{def:module-containment}), which measures the same concept at a coarser granularity: there, the denominator is $|E_{\mathrm{module}}|$ (import edges between modules) and the grouping is by dot-separated module path; here, the denominator is $|E_{\mathrm{thm}}|$ (all declaration-level dependency edges) and the grouping is by namespace truncation. The two metrics are thus not directly comparable in magnitude, since they operate on different edge sets.

A containment ratio of $1$ would mean all dependencies stay within namespace boundaries; $0$ would mean every edge crosses a boundary. The decay is steep: from $48.9\%$ at the top level to ${\sim}14\%$ at depth~$2$, with the steepest drop at the discipline--sub-discipline transition. Full data appear in Appendix~\ref{app:ns-detail} (Table~\ref{tab:containment-decay}, Figure~\ref{fig:containment-curve}).

Degree distribution, centrality, community detection, and robustness analyses for the namespace graph are reported in Appendix~\ref{app:ns-detail}.
